\subsection{The complete initial-data formulation}

Let $(M,g,K)$ be a connected orientable smooth initial data set without
boundary, with $g$ complete.  Suppose that outside a compact set it has
finitely many asymptotically flat ends.  On each end assume
\[
 g-\delta=O_2(r^{-q}),\qquad K=O_1(r^{-1-q}),\qquad q>\tfrac12,
\]
and finite ADM energy and momentum limits as in
Equations~\eqref{n3:intro:energy} and~\eqref{n3:intro:momentum}.
The exponents may differ between ends; their minimum is still larger
than $1/2$.  Assume that the constraint densities
\eqref{n3:intro:constraints} satisfy $\mu\ge |J|_g$ and that $\mu$ and
$|J|_g$ are integrable on $M$.

\begin{definition}[Admissible boundary and minimum enclosing area]
\label{bridge:intro:global-admissible}
Fix an asymptotically flat end of $M$.  An admissible boundary is a
nonempty compact smooth embedded two-sided surface $S=\partial D^+$,
possibly disconnected, where $D^+$ is a smooth open region containing
the entire sufficiently distant part of the chosen end.  The region
$D^+$ need not be connected and may contain other ends.  For the normal
$\nu$ pointing into $D^+$, require
\[
 H_S+\operatorname{tr}_S K\le0.
\]
Define
\begin{equation}
\label{bridge:intro:global-area}
 A_{\min}(S)=\inf\left\{
 \operatorname{Area}_g(\partial D'^+):
 \begin{array}{l}
 D'^+\subset D^+,\quad D'^+\text{ contains the chosen end at infinity},\\
 \partial D'^+\text{ is a compact smooth embedded boundary}
 \end{array}\right\}.
\end{equation}
Competitors may be disconnected or coincide with $S$.  They need not
be trapped or minimal.  No condition on the inward null expansion is
imposed.
\end{definition}

\begin{theorem}[Spacetime Penrose inequality for complete initial data]
\label{bridge:intro:global-penrose}
For the complete initial data just specified, every chosen end and
every admissible boundary satisfy
\begin{equation}
\label{bridge:intro:global-inequality}
 E>|P|_\delta,\qquad
 \sqrt{E^2-|P|_\delta^2}
       \ge \sqrt{\frac{A_{\min}(S)}{16\pi}}.
\end{equation}
Here $(E,P)$ is the ADM vector of the chosen end.  The assertion allows
finitely many ends and disconnected admissible boundaries.
\end{theorem}

\subsection{Removing the timelike hypothesis and recovering the global statement}
\label{bridge:bridge:section}

The timelike exterior inequality already proved also rules out a
non-timelike ADM vector for an exterior with weakly future outer trapped
boundary.  This observation then gives the complete initial-data
formulation by restricting to a suitable one-ended exterior.

\begin{lemma}[Positive enclosing area]
\label{bridge:bridge:positive-area}
Every exterior in Definition~\ref{n3:intro:exterior-class} satisfies
$0<a_g(B)<\infty$, where $B$ is its boundary.
\end{lemma}

\begin{proof}
A sufficiently large coordinate sphere gives finiteness.  For positivity,
choose a small disk in one component of $B$ and a smooth proper embedded
ray from its center to the end, normal to $B$ initially and straight in
the asymptotic coordinates outside a compact set.  The ray has a tubular
neighborhood with disk cross-sections; its coordinate width can be fixed
and positive on the straight tail.  In product coordinates on this tube,
pull back a smooth two-form compactly supported in the transverse disk
and with integral one.  Extend it by zero through the lateral boundary
of the tube.  The resulting two-form $\omega$ is smooth and closed on
the exterior, including its boundary, and has bounded pointwise norm.
The bound follows from compactness on the initial part and asymptotic
flatness on the tail.

Orient the transverse disk so that the flux of $\omega$ through a large
coordinate sphere is one.  A smooth enclosing cut in the interior and a
sufficiently large sphere bound a compact region disjoint from $B$.
Stokes' theorem therefore gives $\int_\Gamma\omega=1$.  The same identity
holds for coincident cuts by smooth outward approximation.  Consequently
\[
 1=\left|\int_\Gamma\omega\right|
 \le \|\omega\|_{L^\infty(g)}\operatorname{Area}_g(\Gamma).
\]
Taking the infimum proves positivity.  If the filled perimeter closure
of the cut class is used, its infimum agrees with the smooth infimum by
Definition~\ref{n3:intro:exterior-class} and the approximation established
in the exterior reduction.
\end{proof}

\begin{corollary}[The exterior inequality without a timelike assumption]
\label{bridge:bridge:timelike}
Let $(N,g,K)$ satisfy all the hypotheses of
Theorem~\ref{n3:intro:exterior-inequality} except the assumption $E>|P|_\delta$.
Then $E>|P|_\delta$ and the conclusion of that theorem holds.
No extension of the data across the compact boundary $B$ is required.
\end{corollary}

\begin{proof}
Use the conformal family of Lemma~\ref{n3:reduction:strict-direction},
which requires no sign assumption on the ADM vector.  For every finite
$s\ge0$, its data $(g_s,K_s)$ satisfy all exterior hypotheses except
possibly timelikeness, with
\[
 E_s=E+sA_\phi,\qquad P_s=P,\qquad A_\phi>0.
\]
Equation~\eqref{n3:reduction:strict-areas} and
Lemma~\ref{bridge:bridge:positive-area} give
$a_{g_s}(B)\ge a_g(B)>0$.
Suppose that $E\le|P|_\delta$ and put
$s_0=(|P|_\delta-E)/A_\phi\ge0$.  For every $s>s_0$,
Equation~\eqref{n3:reduction:strict-charges} places $(N,g_s,K_s)$ in the
timelike class of the already proved
Theorem~\ref{n3:intro:exterior-inequality}.  That theorem and
the preceding area bound imply
\[
 \sqrt{(E+sA_\phi)^2-|P|_\delta^2}
       \ge\sqrt{\frac{a_g(B)}{16\pi}}>0.
\]
Letting $s\downarrow s_0$ gives a contradiction.  Therefore
$E>|P|_\delta$, and applying Theorem~\ref{n3:intro:exterior-inequality}
to the original data proves the claimed inequality.
\end{proof}

\begin{proof}[Proof of Theorem~\ref{bridge:intro:global-penrose}]
Let $D_0$ be the connected component of $D^+$ containing the distant
chosen end.  Its boundary is a union of components of $S$ and is nonempty:
otherwise $D_0$ would be both open and closed in the connected manifold
$M$, hence equal to $M$, contrary to $S\ne\varnothing$.
Put $X_0=\overline{D_0}$, including its smooth compact boundary.
Because $S$ is compact, each sufficiently distant end of $M$ lies
entirely in $D_0$ or entirely outside it.  Outside a compact set, $X_0$
is precisely the union of the original distant ends contained in
$D_0$.  In particular, $X_0$ has finitely many ends, all asymptotically
flat.

For every end of $X_0$ other than the chosen one, choose a sufficiently
distant open coordinate tail, and let $N$ be $X_0$ with all
these open tails removed.  These tails can be chosen disjoint and
outside $S$.  Removing
them preserves connectedness: an excursion of a path into a removed
tail can be replaced by a path on its connected coordinate sphere.
The region $N$ is a smooth manifold with nonempty compact boundary
$B$, consisting of $\partial D_0$ and the added coordinate spheres,
and its only end is the chosen one.

The region is closed in the complete manifold $(M,g)$ and has smooth
boundary.  It is complete for its intrinsic distance with the boundary
included: an intrinsic Cauchy sequence is ambient Cauchy, its ambient
limit lies in $N$, and smooth boundary charts compare the intrinsic
distance locally with Euclidean distance on a half-ball.
The restricted constraints, integrability, and decay satisfy the
exterior hypotheses.  The chosen ADM vector is still $(E,P)$.

On the inherited part of $B$, the normal into $N$ agrees with the normal
into $D^+$, so the assumed expansion is nonpositive.  On an added sphere
in another end, the normal into $N$ points toward decreasing radius.
Consequently
\[
 H_B=-\frac2r+O(r^{-1-q}),\qquad
 \operatorname{tr}_B K=O(r^{-1-q}),
\]
and its future outward expansion is strictly negative when its radius
is sufficiently large.  Corollary~\ref{bridge:bridge:timelike} therefore applies
to $(N,g,K)$ and gives
\begin{equation}
\label{bridge:bridge:restricted-inequality}
 E>|P|_\delta,\qquad
 \sqrt{E^2-|P|_\delta^2}\ge\sqrt{\frac{a_g(B)}{16\pi}}.
\end{equation}

Finally,
\begin{equation}
\label{bridge:bridge:global-area-comparison}
                         a_g(B)\ge A_{\min}(S).
\end{equation}
To see this, first take an enclosing cut $\Gamma$ lying in the interior
of $N$.  The component on its distant chosen-end side is an open region
$D'^+\subset D_0\subset D^+$, with compact smooth boundary consisting
of components of $\Gamma$.  It is a competitor in
Equation~\eqref{bridge:intro:global-area}; its area is at most that of $\Gamma$.
If a cut coincides with any part of $B$, approximate it by smooth cuts
pushed into the interior of $N$, with areas converging to its area.
This also handles cuts coinciding with a newly introduced sphere.
Taking the infimum proves
Equation~\eqref{bridge:bridge:global-area-comparison}.  Other components of
$D^+$ cause no difficulty, since the original competitor class permits
the chosen-end component alone.  Combining
Equations~\eqref{bridge:bridge:restricted-inequality} and
\eqref{bridge:bridge:global-area-comparison} proves the theorem.
\end{proof}

The conformal family above is used only to establish the numerical
inequality and to exclude a non-timelike ADM vector.  The equality
argument is carried out on the original exterior under its separately
stated hypotheses; no equality conclusion is passed through this family
or through the truncation of other ends.
